% GENERATED FILE. Edit canonical lessons, then run npm run generate:books.
% canonical-source-hash: fnv1a64:d53888f5a19bba71
% canonical-lessons: PA-C108-inna, PA-C108-unna, PA-C108-ive, PA-C108-uve, PA-C108-hor

\chapter{This much, That much, Like this, Like that, Other}
\label{ch:pa-this-much-that-much-like-this-like-that-other}
\hlchaptermodality{108}

\begin{chapteropening}
\textbf{By the end of this chapter:} \emph{I can say this much, that much, like this, like that and other.}

Complete the last lesson of chapter 108: i can say this much, that much, like this, like that and other.
\end{chapteropening}

\section[\texorpdfstring{\pa{ਇੰਨਾ} (innā)}{ਇੰਨਾ (innā)}]{\pa{ਇੰਨਾ} (innā) — this much}

\label{lesson:PA-C108-inna}



\begin{quote}
Before the new one: say the Punjabi for both, then the Punjabi for several.
\end{quote}

\subsection*{You'll want to know: \pa{ਇੰਨਾ}}
\textbf{\pa{ਇੰਨਾ}} — \emph{innā} — \textquotedblleft{}this much\textquotedblright{}.

\begin{grammarlens}[title={What you've built}]
One more word for this chapter.
\end{grammarlens}

\subsection*{Your turn}
\begin{itemize}
\raggedright
  \item \emph{Say it:} \emph{innā}
  \item \emph{Say it:} \emph{innā}, once more
  \item \emph{Say it:} say \emph{innā}
  \item \emph{Recall:} say the Punjabi for both, then the Punjabi for several, then say \emph{innā} again
\end{itemize}

\begin{tcolorbox}[breakable,colback=teal!4,colframe=teal!35!black,title={Before you move on}]
What does \pa{ਇੰਨਾ} mean? (\textquotedblleft{}This much\textquotedblright{}.) Say it once more.
\end{tcolorbox}
\section[\texorpdfstring{\pa{ਉੰਨਾ} (unnā)}{ਉੰਨਾ (unnā)}]{\pa{ਉੰਨਾ} (unnā) — that much}

\label{lesson:PA-C108-unna}



\begin{quote}
Before the new one: say the Punjabi for several, then the Punjabi for this much.
\end{quote}

\subsection*{You'll want to know: \pa{ਉੰਨਾ}}
\textbf{\pa{ਉੰਨਾ}} — \emph{unnā} — \textquotedblleft{}that much\textquotedblright{}.

\begin{grammarlens}[title={What you've built}]
One more word for this chapter.
\end{grammarlens}

\subsection*{Your turn}
\begin{itemize}
\raggedright
  \item \emph{Say it:} \emph{unnā}
  \item \emph{Say it:} \emph{unnā}, once more
  \item \emph{Say it:} say \emph{unnā}
  \item \emph{Recall:} say the Punjabi for several, then the Punjabi for this much, then say \emph{unnā} again
\end{itemize}

\begin{tcolorbox}[breakable,colback=teal!4,colframe=teal!35!black,title={Before you move on}]
What does \pa{ਉੰਨਾ} mean? (\textquotedblleft{}That much\textquotedblright{}.) Say it once more.
\end{tcolorbox}
\section[\texorpdfstring{\pa{ਇਵੇਂ} (ivẽ)}{ਇਵੇਂ (ivẽ)}]{\pa{ਇਵੇਂ} (ivẽ) — like this}

\label{lesson:PA-C108-ive}



\begin{quote}
Before the new one: say the Punjabi for this much, then the Punjabi for that much.
\end{quote}

\subsection*{You'll want to know: \pa{ਇਵੇਂ}}
\textbf{\pa{ਇਵੇਂ}} — \emph{ivẽ} — \textquotedblleft{}like this\textquotedblright{}.

\begin{grammarlens}[title={What you've built}]
One more word for this chapter.
\end{grammarlens}

\subsection*{Your turn}
\begin{itemize}
\raggedright
  \item \emph{Say it:} \emph{ivẽ}
  \item \emph{Say it:} \emph{ivẽ}, once more
  \item \emph{Say it:} say \emph{ivẽ}
  \item \emph{Recall:} say the Punjabi for this much, then the Punjabi for that much, then say \emph{ivẽ} again
\end{itemize}

\begin{tcolorbox}[breakable,colback=teal!4,colframe=teal!35!black,title={Before you move on}]
What does \pa{ਇਵੇਂ} mean? (\textquotedblleft{}Like this\textquotedblright{}.) Say it once more.
\end{tcolorbox}
\section[\texorpdfstring{\pa{ਉਵੇਂ} (uvẽ)}{ਉਵੇਂ (uvẽ)}]{\pa{ਉਵੇਂ} (uvẽ) — like that}

\label{lesson:PA-C108-uve}



\begin{quote}
Before the new one: say the Punjabi for that much, then the Punjabi for like this.
\end{quote}

\subsection*{You'll want to know: \pa{ਉਵੇਂ}}
\textbf{\pa{ਉਵੇਂ}} — \emph{uvẽ} — \textquotedblleft{}like that\textquotedblright{}.

\begin{grammarlens}[title={What you've built}]
One more word for this chapter.
\end{grammarlens}

\subsection*{Your turn}
\begin{itemize}
\raggedright
  \item \emph{Say it:} \emph{uvẽ}
  \item \emph{Say it:} \emph{uvẽ}, once more
  \item \emph{Say it:} say \emph{uvẽ}
  \item \emph{Recall:} say the Punjabi for that much, then the Punjabi for like this, then say \emph{uvẽ} again
\end{itemize}

\begin{tcolorbox}[breakable,colback=teal!4,colframe=teal!35!black,title={Before you move on}]
What does \pa{ਉਵੇਂ} mean? (\textquotedblleft{}Like that\textquotedblright{}.) Say it once more.
\end{tcolorbox}
\section[\texorpdfstring{\pa{ਹੋਰ} (hor)}{ਹੋਰ (hor)}]{\pa{ਹੋਰ} (hor) — other, more}

\label{lesson:PA-C108-hor}



\begin{quote}
Before the new one: say the Punjabi for like this, then the Punjabi for like that.

Two Punjabi verbs mean \textquotedblleft{}to see.\textquotedblright{} Which is the more Punjabi-sounding one, and which did Hindi keep? (\emph{Vekhṇā}; Hindi kept \emph{dekhṇā}.)
\end{quote}

\subsection*{You'll want to know: \pa{ਹੋਰ}}
\textbf{\pa{ਹੋਰ}} — \emph{hor} — \textquotedblleft{}other, more\textquotedblright{}.

\begin{grammarlens}[title={What you've built}]
One more word for this chapter.
\end{grammarlens}

\subsection*{Your turn}
\begin{itemize}
\raggedright
  \item \emph{Say it:} \emph{hor}
  \item \emph{Say it:} \emph{hor}, once more
  \item \emph{Say it:} say \emph{hor}
  \item \emph{Recall:} say the Punjabi for like this, then the Punjabi for like that, then say \emph{hor} again
\end{itemize}

\begin{tcolorbox}[breakable,colback=teal!4,colframe=teal!35!black,title={Before you move on}]
What does \pa{ਹੋਰ} mean? (\textquotedblleft{}Other, more\textquotedblright{}.) Say it once more.
\end{tcolorbox}
